% =====================================================================
% sec/2_related.tex  --  Related Work
% =====================================================================

\section{Related Work}
\label{sec:related}

\paragraph{Polyp Segmentation.}
Early CNN-based methods~\cite{fan2020pranet,wei2021shallow} achieved strong
performance on single-dataset benchmarks but struggle to generalize across
colonoscopy centers due to variation in illumination, camera models, and polyp
appearance. Transformer-based architectures~\cite{zhang2021transfuse,dong2021polyp}
improve long-range context modelling but still require domain-specific fine-tuning.

\paragraph{Vision-Language Pre-training for Medical Imaging.}
CLIP~\cite{radford2021clip} pioneered contrastive image-text alignment at scale.
MedCLIP~\cite{wang2022medclip} and UniMedCLIP~\cite{} extend this to the medical
domain using curated radiology and pathology report pairs. MedCLIPSeg~\cite{}
demonstrates that VLP representations can guide segmentation through text prompts,
achieving competitive zero-shot performance on multiple medical modalities.

\paragraph{Parameter-Efficient Fine-Tuning.}
LoRA~\cite{hu2022lora} introduces low-rank weight updates that substantially reduce
the number of trainable parameters. Adapter-based methods~\cite{houlsby2019adapter}
insert lightweight bottleneck modules between frozen Transformer layers.
In this work we apply LoRA to the cross-attention projection matrices inside the
PVL adapters of MedCLIPSeg, combining adapter-style modality bridging with
low-rank efficiency.

\paragraph{Uncertainty Estimation in Medical Segmentation.}
Monte Carlo (MC) dropout~\cite{gal2016dropout} provides a practical approximation
to Bayesian inference. Combining MC-dropout with test-time augmentation (TTA) and
inverse-variance fusion~\cite{wang2019aleatoric} yields well-calibrated uncertainty
maps that are actionable for clinical risk stratification.
